\begin{abstract}
Pathology foundation models readily yield image--label associations, but pooled performance
does not establish transport or information beyond clinical covariates. We tested whether
signals retained support across cohort transport, confounding, site, representation, and
endpoint audits. We implemented a three-state qualification framework---transportable,
context-sensitive, and unsupported in the frozen design---for frozen CONtrastive learning
from Captions for Histopathology (CONCH) and Virchow signals across five prostate-pathology
resources. Targets were grade, phenotype, phosphatase and tensin homolog (PTEN) loss,
speckle type BTB/POZ protein (SPOP) mutation, androgen-receptor (AR) activity, and post-hoc
recurrence. Grade and phenotype showed the strongest cross-cohort support. PTEN remained
detectable in The Cancer Genome Atlas prostate adenocarcinoma (TCGA-PRAD) cohort but lacked
a supported held-out increment beyond grade. AR retained a positive pooled direction while
site transport remained unresolved. SPOP was unsupported in the frozen primary
configuration and setting-sensitive. Recurrence estimates varied by endpoint, encoder, and
scale and lacked established incremental value beyond fuller clinical and site covariates.
Thus, pooled performance separated into evidence states rather than a model ranking. The
framework and empirical failure map identify when apparently successful probes are
transportable, conditional, or unsupported and prioritize signals for further validation.
\end{abstract}
